% ═══════════════════════════════════════════════════════════════
%  PART X — APRIL 8, 2026: THE TANGLED UNIVERSE
% ═══════════════════════════════════════════════════════════════
\clearpage
\part{The Tangled Universe: Polyhedra, Hurwitz Surfaces, and Physical Chirality}

\section{Tangled Platonic Polyhedra and W(3,3)}
\label{sec:tangled}

Hyde and Evans~\cite{HydeEvans2022} construct maximally symmetric
tangled Platonic polyhedra by winding $n$-strand helices on polytori
formed by tubifying the edges of classical polyhedra.
Every number in this construction is a W(3,3) parameter.

\begin{theorem}[Polytorus Genus = W(3,3) Parameters]
Tubifying the edges of each Platonic polyhedron $\{f,z\}$ gives
a polytorus of genus $\mathfrak{g}=1+E-V$:
\begin{center}
\begin{tabular}{@{}lcccc@{}}
\toprule
Polyhedron & $\{f,z\}$ & $V$ & $E$ & $\mathfrak{g}$ \\
\midrule
Tetrahedron & $\{3,3\}$ & 4 & 6 & $\mathbf{3=q}$ \\
Cube & $\{4,3\}$ & 8 & 12 & 5 \\
Octahedron & $\{3,4\}$ & 6 & 12 & $\mathbf{7=\Phisix}$ \\
Icosahedron & $\{3,5\}$ & 12 & 30 & 19 \\
Dodecahedron & $\{5,3\}$ & 20 & 30 & 11 \\
\bottomrule
\end{tabular}
\end{center}
\end{theorem}

The tetrahedron polytorus has genus~$q$,
and the octahedron polytorus has genus~$\Phisix$.
Klein's quartic curve---the prototypical Hurwitz surface---is
\emph{also} genus~$q=3$, establishing a direct link between
tangled tetrahedra and the Klein geometry.

The first non-trivial tangling uses $q=3$ helical strands.
Tangled icosahedra with chiral symmetry~$235=I$ retain
$k=12$ vertices and $2g=30$ edges, and their face cycles
form $(q,\lambda)=(3,2)$ torus knots---trefoils.


\section{The Hurwitz Bound Constant is $k\Phisix$}
\label{sec:hurwitz-bound}

\begin{theorem}[Hurwitz Decomposition]
The Hurwitz bound constant decomposes as:
\begin{equation}
\boxed{84 = k\times\Phisix = 12\times 7.}
\end{equation}
At genus $q=3$ (Klein's quartic):
$|\Aut(\Sigma_q)| = 84(q{-}1) = k\Phisix\lambda = 168 = f\cdot\Phisix$.
The automorphism group is $\PSL(2,\Phisix)=\PSL(2,7)$.
\end{theorem}


\section{The Hurwitz Triplet at Genus $\dim(G_2)$}
\label{sec:hurwitz-triplet}

Bokowski and collaborators~\cite{Bokowski2025} construct polyhedral
embeddings of all Hurwitz surfaces up to genus~14.
The genus-14 case is remarkable: $14 = \lambda\Phisix = \dim(G_2)$.

\begin{theorem}[Hurwitz Triplet Data]
The three genus-14 Hurwitz surfaces each have:
\begin{equation}
V = k\cdot\Phithree = 156, \quad
E = \Phisix\cdot\Phithree\cdot q! = 546, \quad
F = \mu\cdot\Phisix\cdot\Phithree = 364.
\end{equation}
In particular:
\begin{equation}
\frac{F}{\Phisix} = 52 = \dim(F_4) = [3]_3!
\end{equation}
The number of faces divided by $\Phisix$ yields the $q$-factorial at $q=3$.
\end{theorem}


\section{Tetrahedral Chains and the Golden Ratio}
\label{sec:tet-chains}

Akpanya \emph{et al.}~\cite{Akpanya2026} prove the first perfect
closed chain of congruent isosceles tetrahedra has
$N = 14 = \dim(G_2)$ tetrahedra, with an infinite family at
$N = (k{-}1) + k\cdot n$ for $n\in\mathbb{N}_0$.
The edge lengths of these tetrahedra include
$(1+\sqrt{5})/2 = \varphi$, the golden ratio---the quantum
dimension of the Fibonacci anyon from the $\SU(2)_q$ modular category.


\section{Tangling as Physical Mechanism}
\label{sec:tangling-physics}

The synthesis across all four contemporary works reveals:
\begin{itemize}[nosep]
\item \textbf{Tangling requires $q=3$ dimensions} (Lock~2),
\item \textbf{First non-trivial helical tangling uses $q$ strands},
\item \textbf{Chirality breaking} (*$2fz \to 2fz$) in tangling
  mirrors \textbf{parity violation} in the electroweak sector,
\item \textbf{Knot type = particle type} via the Jones polynomial
  and the $\SU(2)_q$ modular functor,
\item \textbf{Hurwitz bound $= k\Phisix$} constrains the maximum
  symmetry of the polytorus hosting the tangle.
\end{itemize}

We conjecture:
\begin{quote}
\emph{The universe is a tangled polytope whose skeleton is the W(3,3)
graph, wound by $q$-strand helices on a polytorus of genus controlled
by the Hurwitz bound $k\Phisix(g{-}1)$. Gauge fields are the helical
windings; chirality is the loss of reflection symmetry upon tangling;
particle generations correspond to surface genus.}
\end{quote}

\begin{thebibliography}{9}
\bibitem{HydeEvans2022} S.\,T. Hyde and M.\,E. Evans,
  ``Symmetric tangled Platonic polyhedra,''
  \emph{PNAS} \textbf{119} (2022) e2110345118.

\bibitem{Bokowski2025} J. Bokowski and K.\,H. (CodeParade),
  ``Polyhedral embeddings of triangular regular maps of genus $g$, $2\leq g\leq 14$,''
  \emph{Symmetry} \textbf{17} (2025) 622.

\bibitem{Akpanya2026} R. Akpanya et al.,
  ``Construction of toroidal polyhedra corresponding to perfect chains of
  isosceles tetrahedra,'' \emph{J.\ Geometry} \textbf{117} (2026) 5.

\bibitem{tangled_nets} S.\,T. Hyde,
  ``Entanglement by design: Symmetry-guided periodic helical windings on crystal nets,''
  arXiv:2603.26817 (2026).
\end{thebibliography}
